% ==============================================================================
% Preamble: Principles of Distributed Quantum Compilers
% Author: Krish Kumar Sharma
% ==============================================================================

\usepackage[utf8]{inputenc}
\usepackage[T1]{fontenc}
\usepackage{lmodern}
\usepackage{microtype}

% --- Page Geometry (Professional Book Layout) ---
\usepackage[
    paper=a4paper,
    top=2.5cm,
    bottom=2.5cm,
    inner=3.0cm,
    outer=2.5cm,
    headheight=26pt,
    footskip=1.2cm
]{geometry}

% --- Mathematics & Quantum Notation ---
\usepackage{amsmath}
\usepackage{amssymb}
\usepackage{amsfonts}
\usepackage{amsthm}
\usepackage{mathtools}
\usepackage{braket} % Dirac notation: \ket{}, \bra{}, \braket{}
\usepackage{bm}
\setcounter{MaxMatrixCols}{30}

% --- Graphics & Visualization ---
\usepackage{graphicx}
\usepackage{xcolor}
\usepackage{tikz}
\usetikzlibrary{shapes,arrows,positioning,fit,backgrounds,calc,matrix,decorations.pathreplacing,decorations.pathmorphing,patterns}

% --- Tables & Typography ---
\usepackage{booktabs}
\usepackage{tabularx}
\usepackage{xltabular}
\usepackage{longtable}
\usepackage{emptypage} % Suppress headers and footers on completely empty pages
\usepackage{array}
\usepackage{multirow}
\usepackage{caption}
\usepackage{subcaption}

% --- Penalty Settings to Eliminate Orphans, Widows, and Isolated Lines ---
\widowpenalty=10000
\clubpenalty=10000
\displaywidowpenalty=10000
\interfootnotelinepenalty=10000
\raggedbottom

% --- Headers & Footers ---
\usepackage{fancyhdr}
\pagestyle{fancy}
\fancyhf{}
\fancyhead[LE]{\bfseries\nouppercase{\leftmark}}
\fancyhead[RO]{\bfseries\nouppercase{\rightmark}}
\fancyfoot[LE,RO]{\thepage}
\renewcommand{\headrulewidth}{0.5pt}
\renewcommand{\footrulewidth}{0pt}

% --- Section Styling ---
\usepackage{titlesec}
\titleformat{\chapter}[display]
  {\normalfont\huge\bfseries\color{darkblue}}
  {\chaptertitlename\ \thechapter}{20pt}{\Huge}

% --- Color Definitions ---
\definecolor{darkblue}{RGB}{16, 44, 87}
\definecolor{quantumviolet}{RGB}{83, 53, 155}
\definecolor{compilerteal}{RGB}{20, 108, 148}
\definecolor{codebg}{RGB}{248, 249, 250}
\definecolor{codeborder}{RGB}{222, 226, 230}
\definecolor{commentgreen}{RGB}{40, 140, 60}
\definecolor{keywordblue}{RGB}{25, 70, 180}

% --- Enhanced List Formatting ---
\usepackage{enumitem}
\setlist[enumerate]{itemsep=4pt, parsep=2pt}
\setlist[itemize]{itemsep=4pt, parsep=2pt}

% --- Professional Charts ---
\usepackage{pgfplots}
\pgfplotsset{compat=1.18}
\usepgfplotslibrary{fillbetween}

% --- Additional Color Definitions for Design Language ---
\definecolor{examplegreen}{RGB}{34, 139, 34}
\definecolor{insightorange}{RGB}{210, 105, 30}
\definecolor{historybrown}{RGB}{139, 90, 43}
\definecolor{cautionred}{RGB}{178, 34, 34}
\definecolor{summaryash}{RGB}{100, 100, 110}
\definecolor{algorithmnavy}{RGB}{20, 40, 100}

% --- Theorem & Definition Environments ---
\usepackage{tcolorbox}
\tcbuselibrary{skins,breakable}

\newtcolorbox{definitionbox}[1]{
    colback=white,
    colframe=compilerteal,
    fonttitle=\bfseries\color{white},
    title=Definition: #1,
    arc=3mm,
    boxrule=1pt,
    breakable
}

\newtcolorbox{theorembox}[1]{
    colback=white,
    colframe=darkblue,
    fonttitle=\bfseries\color{white},
    title={Theorem: #1},
    arc=3mm,
    boxrule=1pt,
    breakable
}

\newtcolorbox{hardwarebox}[1]{
    colback=codebg,
    colframe=quantumviolet,
    fonttitle=\bfseries\color{white},
    title={Hardware Real-World Note: #1},
    arc=3mm,
    boxrule=1pt,
    breakable
}

\newtcolorbox{passbox}[1]{
    colback=codebg,
    colframe=darkblue,
    fonttitle=\bfseries\color{white},
    title={Compiler Pass Mechanics: #1},
    arc=3mm,
    boxrule=1pt,
    breakable
}

% --- NEW DESIGN LANGUAGE ENVIRONMENTS ---

% Worked Examples: Step-by-step walkthroughs with green accents
\newtcolorbox{examplebox}[1]{
    colback=examplegreen!5,
    colframe=examplegreen,
    fonttitle=\bfseries\color{white},
    title={Worked Example: #1},
    arc=3mm,
    boxrule=1.2pt,
    breakable,
    left=6pt,
    right=6pt,
    top=4pt,
    bottom=4pt
}

% Insight Boxes: "Why This Matters" intuition builders
\newtcolorbox{insightbox}[1]{
    colback=insightorange!5,
    colframe=insightorange,
    fonttitle=\bfseries\color{white},
    title={Why This Matters: #1},
    arc=3mm,
    boxrule=1.2pt,
    breakable,
    left=6pt,
    right=6pt
}

% Historical Context: Origin stories and milestones
\newtcolorbox{historicalbox}[1]{
    colback=historybrown!5,
    colframe=historybrown,
    fonttitle=\bfseries\color{white},
    title={Historical Context: #1},
    arc=3mm,
    boxrule=1.2pt,
    breakable,
    left=6pt,
    right=6pt
}

% Caution Boxes: Warnings, common mistakes, pitfalls
\newtcolorbox{cautionbox}[1]{
    colback=cautionred!5,
    colframe=cautionred,
    fonttitle=\bfseries\color{white},
    title={Caution: #1},
    arc=3mm,
    boxrule=1.2pt,
    breakable,
    left=6pt,
    right=6pt
}

% Algorithm Boxes: Pseudocode and formal procedures
\newtcolorbox{algorithmbox}[1]{
    colback=algorithmnavy!3,
    colframe=algorithmnavy,
    fonttitle=\bfseries\color{white},
    title={Algorithm: #1},
    arc=3mm,
    boxrule=1.2pt,
    breakable,
    left=6pt,
    right=6pt
}

% Section Summary Boxes: Recap of key points
\newtcolorbox{summarybox}[1]{
    colback=summaryash!8,
    colframe=summaryash,
    fonttitle=\bfseries\color{white},
    title={Chapter Summary: #1},
    arc=3mm,
    boxrule=1.2pt,
    breakable,
    left=6pt,
    right=6pt
}

% --- Custom Terminology Commands ---
% Highlighted key terms (bold + colored)
\newcommand{\keyterm}[1]{\textbf{\textcolor{compilerteal}{#1}}}

% MLIR keywords in monospace with compiler color
\newcommand{\mlirkw}[1]{\texttt{\textcolor{keywordblue}{#1}}}

% --- Code Listings (MLIR, LLVM IR, C++) ---
\usepackage{listings}
\lstdefinelanguage{MLIR}{
    keywords={module, func, dqc, mpi, llvm, return, attributes},
    keywordstyle=\color{keywordblue}\bfseries,
    ndkeywords={alloc_qubit, h, x, y, z, rx, ry, rz, cnot, cz, swap, ccx, mcx, measure, reset, epr_alloc, telegate, distribute_epr, telegate_sequence},
    ndkeywordstyle=\color{compilerteal}\bfseries,
    identifierstyle=\color{black},
    sensitive=false,
    comment=[l]{//},
    morecomment=[s]{/*}{*/},
    commentstyle=\color{commentgreen}\ttfamily,
    stringstyle=\color{purple}\ttfamily,
    morestring=[b]',
    morestring=[b]"
}

\lstset{
    language=MLIR,
    backgroundcolor=\color{codebg},
    basicstyle=\ttfamily\footnotesize,
    breakatwhitespace=false,
    breaklines=true,
    captionpos=b,
    keepspaces=true,
    numbers=left,
    numbersep=8pt,
    numberstyle=\tiny\color{gray},
    showspaces=false,
    showstringspaces=false,
    showtabs=false,
    tabsize=2,
    frame=single,
    rulecolor=\color{codeborder}
}

% --- Cross-Referencing & Hyperlinks ---
\usepackage{hyperref}
\hypersetup{
    colorlinks=true,
    linkcolor=darkblue,
    citecolor=compilerteal,
    filecolor=magenta,      
    urlcolor=quantumviolet,
    pdfauthor={Krish Kumar Sharma},
    pdftitle={Principles of Distributed Quantum Compilers}
}
\usepackage{cleveref}
